\section*{Chapter \ref{ch:rs:risk-models} --- Risk Models}

\begin{solution}{exo:rs:risk-models:1}
$\sqrt{3.26^2 + (1.23 \times 7.2)^2} = 9.4\%$, against the realised 9.2\%; the \omterm{def:rs:risk-models:bias}{bias statistic} is about $9.4/3.26 = 2.9$ (the measured one is 2.86).
\end{solution}

\begin{solution}{exo:rs:risk-models:2}
$1 \pm \sqrt{2/60} = 1 \pm 0.18$; $1 \pm \sqrt{2/2\,016} = 1 \pm 0.031$.
\end{solution}

\begin{solution}{exo:rs:risk-models:3}
The industry dummies sum to the country factor's exposure (one for every stock), so the regression is singular: any constant can move between the country
factor and the industries. Constraining the cap-weighted industry returns to zero makes the country factor the cap-weighted market and each \omterm{def:rs:risk-models:style}{industry factor}
the industry's return net of the market (a dollar-neutral industry portfolio).
\end{solution}

\begin{solution}{exo:rs:risk-models:4}
Each forecast is the true variance plus estimation noise. Sorting by the forecast sorts partly by the noise: the top decile holds the stocks whose recent residuals
were luckily large (true variance below the forecast), the bottom those luckily small. The realised variance regresses to the true one, so the top decile is
overpredicted (0.93) and the bottom underpredicted (1.13). Shrinking towards a group mean undoes part of the sort.
\end{solution}

\begin{solution}{exo:rs:risk-models:5}
The adjustment scales the factor covariance by the recent cross-sectional bias of the factors. The market portfolio's risk is almost all factor risk, whose regimes
the adjustment tracks. The random books' risk is mostly specific; they gain nothing, and inherit the multiplier's noise on their small factor part.
\end{solution}

\begin{solution}{exo:rs:risk-models:6}
$\sum_k x_k (Fx)_k = x^\top F x$. A contribution is negative when the exposure to a factor hedges the others: $x_k$ and $(Fx)_k$ of opposite signs.
\end{solution}

\begin{solution}{exo:rs:risk-models:7}
Run \texttt{rs\_riskmodel.summary(0.3)}. The momentum exposure falls, so the model without momentum underpredicts less: \omterm{def:rs:risk-models:bias}{bias statistics} of 1.98
without momentum, 1.02 with all factors and 1.48 for the statistical model (realised volatility 6.53\%, forecast 3.34\% without momentum).
\end{solution}

\begin{solution}{exo:rs:risk-models:8}
The $R^2$ is about single stocks, whose daily variance is mostly specific. A book's risk is the factor part, which does not diversify, plus a specific part that
falls with the number of names; the model's value is in forecasting both, measured by \omterm{def:rs:risk-models:bias}{bias statistics} (1.04 on the exposed book, 1.00 on average for random books),
not by the regression's fit.
\end{solution}

\begin{solution}{pb:rs:risk-models:1}
\begin{enumerate}
\item Country (one for every stock); ten industries (membership); beta (trailing-year regression on the cap-weighted market, monthly); size, value, momentum (size fixed at listing; value from the latest filed book; momentum from the 12--1 month return; all known at the previous close).
\item Specific variance falls with capitalisation; the weights approximate generalised least squares and give the country factor the cap-weighted market's meaning.
\item 17.8\%; between 8.0\% and 8.7\%; 7.2\%.
\item 0.19: a stock's daily variance is mostly specific, so \omterm{def:rs:risk-models:specific}{specific risk} matters even in a book of two hundred names.
\item It multiplies the factor covariance by an exponentially weighted average of the cross-sectional bias; the market's rolling one-year \omterm{def:rs:risk-models:bias}{bias statistics} range
over 0.56--1.32 instead of 0.50--1.55, and 44\% of windows fall outside the band instead of 66\%.
\item Their risk is mostly specific; the random books' share inside the band is 92\% with it and 98\% without.
\item 1.13 in the lowest decile to 0.93 in the highest; shrinking 10\% towards the size decile's mean gives 0.97 to 1.04.
\item The Student-$t$ shocks and earnings-day jumps of the simulation, which an exponentially weighted variance reflects late.
\item A score of 0.5 times the momentum score plus $\sqrt{0.75}$ times noise; the top and bottom hundred names; neutralised to the model without momentum;
gross 2; rebuilt monthly.
\item All factors: 9.23\% forecast, 9.20\% realised, 1.04. Without momentum: 3.26\%, 9.20\%, 2.86. Statistical: 5.59\%, 9.20\%, 1.68.
\item 76\% of the forecast variance on day 1\,500, at an exposure of 1.03 (root mean square 1.23 over the eight years).
\item Momentum is a direction of common variation, which principal components find; but its loadings are reset monthly, and a trailing year's components
assume loadings that stay put.
\item \textbf{Named result.} \omterm{def:rs:risk-models:bias}{Bias statistic} 2.86 for the model missing momentum; realised volatility 2.83 times the forecast (9.20\% against 3.26\%).
\item By running a statistical model beside the fundamental one (5.59\% against 3.26\%), or by regressing the book's returns on candidate \omterm{def:rs:risk-models:families}{factor returns}.
\item \omterm{def:rs:risk-models:bias}{Bias statistics} of random books, factor portfolios and the firm's books, rolling and over the whole sample, compared with the incumbent model.
\item The band is too narrow: correct forecasts fall outside it more than 5\% of the time (86\% inside at a kurtosis of 5 over 120 periods, in MSCI's figures).
\item A fundamental model with statistical factors added (for example, principal components of its residuals); worth it when books take risks the characteristics
do not describe.
\item Random books average over exposures; a model can be right on average and wrong for the particular exposure a real book concentrates.
\item The book the firm holds, and books built from its signals, because a signal can load on a factor the model lacks.
\item The risks it has names for, and nothing else.
\end{enumerate}
\end{solution}

\begin{solution}{iq:rs:risk-models:1}
Fundamental: observed exposures, \omterm{def:rs:risk-models:families}{factor returns} by cross-sectional regression; interpretable, stable, good for attribution and constraints. Statistical: both
estimated from returns; finds unnamed common variation and adapts, but its factors are hard to interpret and unstable. Prefer fundamental for construction and
reporting, statistical as a check and for short horizons; hybrids combine them.
\end{solution}

\begin{solution}{iq:rs:risk-models:2}
Each day, weighted least squares of the cross-section of returns on the previous close's exposures, weights the square root of capitalisation, with the
cap-weighted industry returns constrained to zero because country and industry dummies are collinear; the constraint makes the country factor the market and the
industries net of it.
\end{solution}

\begin{solution}{iq:rs:risk-models:3}
Decompose the realised P\&L onto the model's factors and the residual; look for a common component in the residuals (principal components of the book's names'
residuals, or regression on candidate \omterm{def:rs:risk-models:families}{factor returns}); compare with a statistical model's forecast; check the regime (a volatility shift the model lags).
\end{solution}

\begin{solution}{iq:rs:risk-models:4}
Standardise many portfolios' returns by their forecasts and compute the standard deviation (the \omterm{def:rs:risk-models:bias}{bias statistic}), overall and in rolling windows, with its band
$1 \pm \sqrt{2/T}$. Pitfalls: fat tails widen the true band; one portfolio is weak evidence; averaging over portfolios hides errors on particular exposures.
\end{solution}

\begin{solution}{iq:rs:risk-models:5}
The optimiser seeks the directions the estimated covariance says are least risky, which are disproportionately those where sampling error made the estimate too
low; Menchero, Wang and Orr correct the eigenportfolios' biases.
\end{solution}

\begin{solution}{iq:rs:risk-models:6}
Inputs: point-in-time prices, returns, capitalisations, industry classifications and descriptors. Order: exposures at the close; the day's factor regression and
residuals; covariance, regime multiplier, specific variances; forecasts. Checks: constraint satisfied, $R^2$ and \omterm{def:rs:risk-models:families}{factor returns} within ranges, no exposure jumps
without cause, \omterm{def:rs:risk-models:bias}{bias statistics} monitored. Store every day's exposures, \omterm{def:rs:risk-models:families}{factor returns}, residuals and forecasts, versioned, so any past forecast can be reproduced.
\end{solution}
